\subsection{Quantitative separation from the same finite data}
\label{sec:fd-arithmetic-separation}

The arithmetic certificate also gives a separation bound on the entire
certified packet. Let \(d_N=\dim V_N\), and let
\(\alpha_N>0\) be the minimum Rayleigh quotient of \(q\) against
the \(L^2\) norm on \(V_N\). Thus
\[
 q(f,f)\ge\alpha_N\|f\|_{L^2(-N,N)}^2\qquad(f\in V_N).
\]
Such a finite bound is supplied by the Gram matrix: if \(M_N\) is the
\(L^2\) Gram matrix in the hat-difference basis, then
\[
 \alpha_N\ge
 \frac{1}{\bigl(\sum_{i,j}|(G^{-1})_{ij}|\bigr)
                 \operatorname{tr}M_N},\qquad
 \operatorname{tr}M_N=\frac{(4d_N-2)2^{-N}}3.
\]
The right-hand side therefore gives an admissible finite lower bound.
Every generic hat difference has squared \(L^2\) norm
\(4\cdot2^{-N}/3\).
The two hats adjacent to the central hat each reduce that value by
\(2^{-N}/3\), proving the trace formula.

\begin{proposition}[A finite separation inequality]
\label{prop:fd-arithmetic-separation}
For every \(N,k\ge1\), \(T>0\), and every two distinct occurrences
in the zero multiset with \(0<\Im\rho\le T\), their
ordinates satisfy
\[
 |\gamma_1-\gamma_2|^2\ge
 \frac{3\alpha_N}{N^3}
       \left(1-\frac{D_{N,k}T^{2k}}2\right).
 \label{eq:fd-separation-bound}
\]
The finite arithmetic data above give the strict bounds
\[
 \alpha_6>7\cdot10^{-17},\qquad
 |\gamma_1-\gamma_2|>6\cdot10^{-10}
       \quad(0<\gamma_1,\gamma_2\le90),
\]
where the two zeros in the second statement are two distinct
occurrences in the multiset, so this bound also excludes a repeated ordinate.
For the reciprocal spectral values, define
\[
 d(T)=\inf\left\{
  \left|\frac1\gamma-\frac1{\gamma'}\right|:
  0<|\gamma|\le T,\ \gamma'\in\Gamma,\ \gamma'\ne\gamma
 \right\}.
\]
Then the same calculation gives \(d(45)>1.48\cdot10^{-13}\).
\end{proposition}

\begin{proof}
Let \(e_\gamma\) be the unit vector of the coordinate indexed by
the distinct signed ordinate \(\gamma\), and put
\(p_N(\gamma)=\|P_Ne_\gamma\|_q^2\).
Choose an orthonormal eigenbasis \(u_j\) of the finite compression.
Scalar convexity of \(x\mapsto x^{2k}\), applied to the spectral
measure of \(C\) in each \(u_j\), gives
\[
 \langle Cu_j,u_j\rangle_q^{2k}
       \le\langle C^{2k}u_j,u_j\rangle_q.
\]
Adding these inequalities and expanding the positive trace in the
coordinate basis proves
\[
 D_{N,k}\ge
 2\sum_{\gamma>0}\frac{m_\gamma-1}{\gamma^{2k}}
 +\sum_{\gamma\in\Gamma}
             \frac{1-p_N(\gamma)}{\gamma^{2k}}.
 \tag{*}
\]
All summands are nonnegative. If a zero in the packet is multiple,
the first sum already gives \(D_{N,k}\ge2/T^{2k}\), and the
asserted separation inequality has a nonpositive right-hand side.
It remains to consider two distinct simple positive ordinates.

Put \(\delta=|\gamma_1-\gamma_2|\) and
\(v_+=(e_{\gamma_1}-e_{\gamma_2})/\sqrt2\).
For \(f\in V_N\), its Fourier transform obeys
\[
 \begin{split}
 |\widehat f(\gamma_1)-\widehat f(\gamma_2)|
 &\le\delta\int_{-N}^N|x f(x)|\,dx\\
 &\le\delta\left(\frac{2N^3}3\right)^{1/2}\|f\|_2.
 \end{split}
\]
The coordinates at these two ordinates have multiplicity one.
Dividing the squared estimate by two, and taking the supremum over
\(q(f,f)=1\), yields
\[
 \|P_Nv_+\|_q^2\le\frac{\delta^2N^3}{3\alpha_N}.
\]
The identical estimate holds for the negative contrast
\(v_-=(e_{-\gamma_1}-e_{-\gamma_2})/\sqrt2\).
In each two-dimensional coordinate space, extend the contrast
to an orthonormal basis. Positivity of \(I-P_N\) then shows
\[
 \sum_{\gamma\in\{\pm\gamma_1,\pm\gamma_2\}}
      (1-p_N(\gamma))
 \ge2\left(1-\frac{\delta^2N^3}{3\alpha_N}\right).
\]
Each of these four ordinates has modulus at most \(T\).
Substitution into \((*)\) and rearrangement proves the inequality.

The Gram lower bound follows from
\(\lambda_{\min}(G)\ge(\sum_{i,j}|(G^{-1})_{ij}|)^{-1}\)
and \(\|M_N\|\le\operatorname{tr}M_N\).
The stored interval bound for the inverse matrix at level six
certifies that the resulting ratio exceeds \(7\cdot10^{-17}\).
Use the already verified rational upper bound
\(D_{6,3}<2.309270917\cdot10^{-12}\).
The remaining numerical comparison is the exact rational inequality
\[
 (6\cdot10^{-10})^2<
 \frac{3(7\cdot10^{-17})}{6^3}
 \left(1-\frac{(2.309270917\cdot10^{-12})90^6}{2}\right).
\]
It proves the claimed separation without a new matrix computation.

Finally consider a signed ordinate \(\gamma\) with
\(0<|\gamma|\le45\). An ordinate of the same sign and modulus
at most \(90\) has reciprocal distance greater than
\((6\cdot10^{-10})/(45\cdot90)\).
An ordinate of the same sign and modulus greater than \(90\)
has reciprocal distance greater than \(1/90\).
An opposite-sign ordinate has reciprocal distance greater than
\(1/45\). Taking the minimum proves
\[
 d(45)\ge\frac{6\cdot10^{-10}}{4050}
          >1.48\cdot10^{-13}.
\]
The \href{\workbenchurl}{\emph{Arithmetic separation} certificate}
records these scalar checks and the hash of their input interval data.
\end{proof}
